% Part: model-theory
% Chapter: interpolation
% Section: interpolation-proof

\documentclass[../../../include/open-logic-section]{subfiles}

\begin{document}

\olfileid{mod}{int}{prf}

\olsection{Teorema Interpolasi Craig}

\begin{thm}[Teorema Interpolasi Craig]
\ollabel{thm:interpol}
Ing andharan iki A lan B ukara, miturut cathetan cakupan SOL6-F026.
Yen $\Entails !A \lif !B$, ana !!a{sentence} $!C$ kang
$\Entails !A \lif !C$ lan $\Entails !C \lif !B$, lan saben
!!{constant}, !!{function}, lan !!{predicate} (saliyane $\eq$) ing
$!C$ uga katon ing $!A$ lan~$!B$. !!{sentence} $!C$ iki diarani
\emph{interpolan} kanggo $!A$ lan~$!B$.
\end{thm}

\begin{proof}
Upamane $\Lang{L}_1$ iku basa saka $!A$ lan $\Lang{L}_2$ iku
basa saka $!B$. Tetepna $\Lang{L}_0 = \Lang{L}_1 \cap \Lang{L}_2$.
Kanggo saben $i \in \{0, 1, 2 \}$, wangun $\Lang{L}'_i$ dipikolehi
saka $\Lang{L}_i$ kanthi nambah tanpa wates akehe !!{constant}s anyar
$\Obj c_0, \Obj c_1, \Obj c_2, \dots$.

Yen $!A$ ora duwe model, $\lexists[x][\eq/[x][x]]$ iku
interpolan. Yen $\lnot !B$ ora duwe model (mula $!B$ valid),
$\lexists[x][\eq[x][x]]$ iku interpolan. Mula kita uga bisa
nganggep $!A$ lan $\lnot !B$ loro-lorone duwe model.

Kanggo mbuktekake kontraposisi Teorema Interpolasi,
upamane ora ana interpolan kanggo $!A$ lan $!B$.
Tegese, upamane $\{!A\}$ lan $\{\lnot !B\}$ ora bisa dipisahake
ing $\Lang{L}_0$.

Tujuan kita yaiku nggedhekake pasangan $(\{ !A \}, \{\lnot!B\})$
dadi pasangan kang ora bisa dipisahake lan maksimal
$(\Gamma^*, \Delta^*)$. Upamane $!A_0$,
$!A_1$, $!A_2$, \dots ndhaptar !!{sentence}s saka $\Lang{L}'_1$,
lan $!B_0$, $!B_1$, $!B_2$, \dots ndhaptar !!{sentence}s
saka~$\Lang{L}'_2$. Cathetan koreksi sumber OLPL-115:
dhaptar kudu nyakup basa kang wis ditambahi konstanta anyar,
amarga maksimalitas ing ngisor iki ditegesake ing basa iku.
Kita tegesi rong barisan mundhak saka himpunan
!!{sentence}s $(\Gamma_n, \Delta_n)$ kanggo $n \ge 0$ kaya iki.
Tetepna $\Gamma_0 = \{ !A\}$ lan $\Delta_0 = \{\lnot !B \}$.
Yen $(\Gamma_n, \Delta_n)$ wis ditegesi, tegese
$\Gamma_{n+1}$ lan $\Delta_{n+1}$ mangkene. Saben himpunan ing tahap anyar nggawa kabeh ukara saka tahap sadurunge; yen calon ukara ditolak, sisih iku tetep. Konstanta kanggo saksi dijupuk saka dhaptar konstanta anyar umum kang ditambahake marang rong basa, supaya durung ana ing himpunan winates ing tahap iku:
\begin{enumerate}
\item Yen $\Gamma_n \cup \{!A_n \}$ lan $\Delta_n$ ora bisa dipisahake
  ing $\Lang{L}'_0$, lebokake $!A_n$ menyang $\Gamma_{n+1}$.
  Kajaba iku, yen $!A_n$ iku !!{formula} eksistensial
  $\lexists[x][!S]$, pilih !!{constant} anyar $c$ kang ora katon
  ing $\Gamma_n$, $\Delta_n$, $!A_n$ utawa $!B_n$, lan lebokake
  $\Subst{!S}{c}{x}$ menyang $\Gamma_{n+1}$.
\item Yen $\Gamma_{n+1}$ lan $\Delta_n \cup \{!B_n \}$ ora bisa
  dipisahake ing $\Lang{L}'_0$, lebokake $!B_n$ menyang
  $\Delta_{n+1}$. Kajaba iku, yen $!B_n$ iku !!{formula}
  eksistensial $\lexists[x][!S]$, pilih !!{constant} anyar $c$
  kang ora katon ing $\Gamma_{n+1}$, $\Delta_n$, $!A_n$
  utawa $!B_n$, lan lebokake $\Subst{!S}{c}{x}$ menyang
  $\Delta_{n+1}$.
\end{enumerate}
Pungkasane, tegesi:
\begin{align*}
  \Gamma^* & = \bigcup_{n\ge 0} \Gamma_n, &
  \Delta^* & = \bigcup_{n\ge 0} \Delta_n.
\end{align*}
Kanthi induksi bebarengan marang $n$, saiki kita bisa mbuktekake:
\begin{enumerate}
\item\ollabel{part-a} $\Gamma_n$ lan $\Delta_n$ ora bisa dipisahake
  ing $\Lang{L}'_0$;
\item\ollabel{part-b} $\Gamma_{n+1}$ lan $\Delta_n$ ora bisa
  dipisahake ing $\Lang{L}'_0$.
\end{enumerate}
Dhasar kanggo \olref{part-a} diwenehake dening
\olref[sep]{lem:sep1}. Kanggo \olref{part-b}, kita kudu
mbedakake telung kasus:
\begin{enumerate}
\item Yen $\Gamma_0 \cup \{!A_0 \}$ lan $\Delta_0$ bisa dipisahake,
  mula $\Gamma_1 = \Gamma_0$ lan \olref{part-b} padha karo
  \olref{part-a};
\item Yen $\Gamma_1 = \Gamma_0 \cup\{ !A_0\}$, mula $\Gamma_1$
  lan $\Delta_0$ ora bisa dipisahake miturut konstruksi.
\item Isih ana kasus nalika $!A_0$ eksistensial, saengga
  $\Gamma_1 = \Gamma_0 \cup \{ \lexists[x][!S], \Subst{!S}{c}{x}
  \}$. Miturut konstruksi, $\Gamma_0 \cup \{ \lexists[x][!S]\}$
  lan $\Delta_0$ ora bisa dipisahake; mula miturut
  \olref[sep]{lem:sep2}, $\Gamma_0 \cup \{ \lexists[x][!S],
  \Subst{!S}{c}{x} \}$ lan $\Delta_0$ uga ora bisa dipisahake.
\end{enumerate}
Iki ngrampungake dhasar induksi kanggo \olref{part-a} lan
\olref{part-b} ing ndhuwur. Saiki menyang langkah induksi.
Kanggo \olref{part-a}, yen $\Delta_{n+1} = \Delta_n \cup
\{ !B_n \}$, mula $\Gamma_{n+1}$ lan $\Delta_{n+1}$ ora bisa
dipisahake miturut konstruksi (sanajan $!B_n$ eksistensial,
miturut \olref[sep]{lem:sep2}); yen $\Delta_{n+1} = \Delta_n$
(amarga $\Gamma_{n+1}$ lan $\Delta_n \cup \{!B_n\}$ bisa
dipisahake), kita gunakake hipotesis induksi \olref{part-b}.
Kanggo langkah induksi \olref{part-b}, yen
$\Gamma_{n+2} = \Gamma_{n+1} \cup \{!A_{n+1} \}$, mula
$\Gamma_{n+2}$ lan $\Delta_{n+1}$ ora bisa dipisahake
miturut konstruksi (sanajan $!A_{n+1}$ eksistensial,
miturut \olref[sep]{lem:sep2}); yen
$\Gamma_{n+2} = \Gamma_{n+1}$, kita gunakake kasus induksi
kanggo \olref{part-a} kang lagi dibuktekake. Iki ngrampungake
induksi kanggo \olref{part-a} lan \olref{part-b}.

Mula $\Gamma^*$ lan $\Delta^*$ ora bisa dipisahake; yen bisa,
miturut kompaktitas ana indeks $n \ge 0$ lan ukara pamisah kang misahake
$\Gamma_n$ lan $\Delta_n$, nalisir \olref{part-a}.
Mligine, $\Gamma^*$ lan $\Delta^*$ konsisten:
yen sing kapisan utawa kapindho ora konsisten,
loro-lorone bisa dipisahake dening
$\lexists[x][\eq/[x][x]]$ utawa
$\lforall[x][\eq[x][x]]$, manut kasusé.

Saiki kita tuduhake yen $\Gamma^*$ konsisten maksimal ing
$\Lang{L}'_1$, lan semono uga $\Delta^*$ ing $\Lang{L}'_2$.
Kanggo sing kapisan, upamane $!A_n \notin \Gamma^*$ lan
$\lnot !A_n \notin \Gamma^*$ kanggo sawenehing $n \ge 0$.
Yen $!A_n \notin \Gamma^*$, mula $\Gamma_n \cup \{!A_n \}$
bisa dipisahake saka $\Delta_n$, dadi ana $!C \in
\Lang{L}'_0$ saengga loro pratelan iki bener:
\begin{align*}
  \Gamma^* & \Entails !A_n \lif !C, &
  \Delta^* & \Entails \lnot !C.
\end{align*}
Semono uga, yen $\lnot !A_n \notin \Gamma^*$, ana $!C' \in
\Lang{L}'_0$ saengga loro pratelan iki bener:
\begin{align*}
  \Gamma^* & \Entails \lnot !A_n \lif !C', &
  \Delta^* & \Entails \lnot !C'.
\end{align*}
Miturut logika proposisional, $\Gamma^* \Entails !C \lor !C'$
lan $\Delta^* \Entails \lnot (!C \lor !C')$, mula $!C \lor
!C'$ misahake $\Gamma^*$ lan $\Delta^*$. Argumen kang padha
nuduhake yen $\Delta^*$ maksimal.

Pungkasan, kita tuduhake yen $\Gamma^* \cap \Delta^*$
konsisten maksimal ing $\Lang{L}'_0$. Cetha yen konsisten,
amarga minangka irisan rong himpunan konsisten.
Kanggo nuduhake maksimalitas, pilih $!S \in \Lang{L}'_0$.
Saiki, $\Gamma^*$ maksimal ing $\Lang{L'_1}
\supseteq \Lang{L'_0}$, lan semono uga $\Delta^*$
maksimal ing $\Lang{L'_2} \supseteq \Lang{L'_0}$.
Mula $!S \in \Gamma^*$ utawa $\lnot !S \in \Gamma^*$,
lan $!S \in \Delta^*$ utawa $\lnot !S \in \Delta^*$.
Yen $!S \in \Gamma^*$ lan $\lnot !S \in \Delta^*$,
$!S$ bakal misahake $\Gamma^*$ lan $\Delta^*$;
yen $\lnot !S \in \Gamma^*$ lan $!S \in \Delta^*$,
$\Gamma^*$ lan $\Delta^*$ bakal dipisahake
dening $\lnot !S$.
Mula $!S \in \Gamma^* \cap \Delta^*$ utawa
$\lnot !S \in \Gamma^* \cap \Delta^*$,
lan $\Gamma^* \cap \Delta^*$ maksimal.

Cathetan pangganep bukti SOL6-F028: maksimalitas wae durung njamin model kang kabeh unsure dijenengi konstanta. Konstruksi iki uga menehi saksi kanggo saben ukara eksistensial kang ana ing himpunan lintang. Yen ukara kasebut ditolak nalika giliran enumerasine, separator ing tahap iku bakal tetep misahake himpunan pungkasan, nalisir inseparability. Mula ukara kasebut ditampa lan diwenehi saksi. Kanggo saben term tertutup t, ukara eksistensial x padha karo t nduweni saksi saka konstanta anyar umum, dadi saben unsur model term bisa dijenengi konstanta saka reservoir kang padha ing rong basa.
Amarga $\Gamma^*$ konsisten maksimal lan nduweni sipat saksi iki, himpunan iki nduweni
model $\Struct{M}'_1$ kang !!{domain} $\Domain{M'_1}$
dumadi saka, lan mung saka, unsur-unsur $\Assign{c}{M'_1}$
kang dadi interpretasi !!{constant}s---kaya ing bukti
teorema kelengkapan
(\olref[fol][com][cth]{thm:completeness}).
Semono uga, $\Delta^*$ nduweni model $\Struct{M}'_2$
kang !!{domain} $\Domain{M'_2}$ diwenehi dening
interpretasi $\Assign{c}{M'_2}$ saka !!{constant}s.

Wangun $\Struct{M_1}$ saka $\Struct{M'_1}$ kanthi
ngilangi interpretasi !!{constant}s, !!{function}s,
lan !!{predicate}s ing $\Lang{L'_1} \setminus
\Lang{L'_0}$; tumrap $\Struct{M_2}$ uga mangkono.
Pemetaan $h \colon M_1 \to M_2$ kang ditegesi dening
$h(\Assign{c}{M'_1}) = \Assign{c}{M'_2}$ iku
isomorfisme ing $\Lang{L}'_0$, amarga
$\Gamma^* \cap \Delta^*$ konsisten maksimal ing
$\Lang{L}'_0$, kaya kang wis dituduhake.
Sebabé, saben $\Lang{L}'_0$-!!{sentence}
kalebu ing $\Gamma^*$ lan $\Delta^*$ loro-lorone,
utawa ora kalebu ing loro-lorone. Mula
$\Assign{c}{M'_1} \in \Assign{P}{M'_1}$ yen lan
mung yen $\Atom{P}{c} \in \Gamma^*$, yen lan mung yen
$\Atom{P}{c} \in \Delta^*$, yen lan mung yen
$\Assign{c}{M'_2} \in \Assign{P}{M'_2}$.
Syarat isomorfisme liyane bisa dibuktekake kanthi cara padha. Ing kene c tansah dijupuk saka konstanta anyar umum. Ukara kesamaan antarane rong konstanta umum padha ing rong himpunan lintang; mula pemetaan iki ora gumantung marang pilihan jeneng lan injektif. Saben unsur model kapindho uga duwe jeneng umum, mula pemetaan iki surjektif.

Saiki tegesi model $\Struct{M}$ kanggo !!{language}
$\Lang{L}'_1 \cup \Lang{L}'_2$ kaya mangkene:
Cathetan koreksi sumber OLPL-117:
basa gabungan iki kudu ngemot konstanta anyar,
supaya pratelan pungkasan ngenani
gabungan rong himpunan lintang bisa ditegesi.
\begin{enumerate}
\item !!{domain} $\Domain{M}$ yaiku $\Domain{M_2}$,
  tegese himpunan kabeh unsur $\Assign{c}{M'_2}$;
\item Yen !!a{predicate}~$P$ ana ing $\Lang{L_2} \setminus
  \Lang{L_1}$, mula $\Assign{P}{M} = \Assign{P}{M'_2}$;
\item Yen predikat $P$ ana ing $\Lang{L}_1\setminus \Lang{L}_2$,
  mula $\Assign{P}{M} = h(\Assign{P}{M'_1})$.
  Cathetan koreksi sumber OLPL-116:
  predikat iki mung ana ing basa kapisan, dadi relasine
  dijupuk saka model kapisan banjur ditranspor nganggo h.
  Tegese, $\tuple{\Assign{c_1}{M'_2}, \dots,
  \Assign{c_n}{M'_2}} \in \Assign{P}{M}$ yen lan mung yen
  $\tuple{\Assign{c_1}{M'_1}, \dots,
  \Assign{c_n}{M'_1}} \in \Assign{P}{M'_1}$.
\item Yen !!a{predicate} $P$ ana ing $\Lang{L}_0$,
  mula $\Assign{P}{M} = \Assign{P}{M'_2} =
  h(\Assign{P}{M'_1})$.
\item !!^{function}s saka $\Lang{L}_1 \cup \Lang{L}_2$,
  kalebu !!{constant}s, ditangani kanthi cara padha. Konstanta anyar umum diwenehi interpretasi saka model kapindho, kang padha karo transport interpretasi model kapisan liwat h.
\end{enumerate}

Pungkasan, induksi marang !!{formula}s nuduhake yen
$\Struct{M}$ cocog karo $\Struct{M'_1}$ kanggo kabeh
!!{formula}s saka $\Lang{L'_1}$ liwat isomorfisme h,
lan cocog karo $\Struct{M'_2}$ kanggo kabeh
!!{formula}s saka $\Lang{L'_2}$.
Cathetan koreksi sumber OLPL-118:
amarga ranah model kapisan lan model gabungan beda,
kecocokan tumrap pangaji variabel kudu liwat h.
Mligine, $\Struct{M} \Entails \Gamma^* \cup \Delta^*$,
mula $\Struct{M} \Entails !A$ lan
$\Struct{M} \Entails \lnot!B$, saengga
$\not\Entails !A \lif !B$. Iki ngrampungake bukti
Teorema Interpolasi Craig.
\end{proof}

\end{document}
